\subsection{与分次模相伴的素理想}

命题 1. 设 \(\Delta\) 是无挠交换群，A 是 \(\Delta\) 型分次环，M 是 \(\Delta\) 型分次 A-模。每个与 M 相伴的素理想都是分次的，且是 M 中某个齐次元的湮灭子。

我们知道，\(\Delta\) 可以赋予一个与其群结构相容的全序结构（代数，第二章，第 11 节，第 4 节，引理 2）。设 p 是与 M 相伴的素理想，是元素 \(x \in M\) 的湮灭子；令 \((x_{i})_{i \in \Delta}\) 为 x 的齐次分量族，并令 \(i(1) < i(2) < \cdots < i(r)\) 为使 \(x_{i} \neq 0\) 的 i 的值。取 \(a \in p\)，令 \((a_{i})_{i \in \Delta}\) 为其齐次分量族；我们将证明对于所有 \(a_{i} \in p\)，都有 \(i \in \Delta\)，这将表明 p 是分次理想。

  我们对指标 \(i\)（满足 \(a_{i} \neq 0\)）的数目作归纳。若该数目为 0，断言显然成立；否则，令 \(m\) 为指标 \(i\) 中满足 \(a_{i} \neq 0\) 的最大者；若能证明 \(a_{m} \in \mathfrak{p}\)，则对 \(a - a_{m}\) 应用归纳假设即可得到结论。现在 \(ax = 0\)；对于所有 \(j \in \Delta\)，利用次数为 \(m + j\) 的 \(ax\) 的齐次分量为 0 这一事实，得到 \(\sum_{i \in \Delta} a_{m - i} x_{j + i} = 0\)；我们从而得到 \(a_{m} x_{j}\) 是 \(x_{i}\) 的线性组合，其指标满足 \(i > j\)。特别地，\(a_{m} x_{i(r)} = 0\)；于是通过对 \(n < r\) 作递降归纳，得到 \(a_{m}^{r - n + 1} x_{i(n)} = 0\)。然后 \(a_{m}^{r} x = 0\)，从而 \(a_{m}^{r} \in \mathfrak{p}\)；由于 \(\mathfrak{p}\) 是素理想，\(a_{m} \in \mathfrak{p}\)。

  现在证明 \(\mathfrak{p}\) 是 M 中某个齐次元的湮灭子。记 \(\mathfrak{b}_n = \operatorname{Ann}(x_{i(n)})\)，其中 \(1 \leqslant n \leqslant r\)。对于 \(b\) 的每个齐次元 \(\mathfrak{p}\) 及所有 \(n\)，\(bx\) 的次数为 \(i(n) + \deg(b)\) 的齐次分量是 \(bx_{i(n)}\)，所以 \(bx_{i(n)} = 0\)，因而 \(b \in \mathfrak{b}_n\)；由于 \(\mathfrak{p}\) 由其齐次元生成，\(\mathfrak{p} \subset \mathfrak{b}_n\)。另一方面，显然 \(\bigcap_{n=1}^{r} \mathfrak{b}_n \subset \operatorname{Ann}(x) = \mathfrak{p}\)；由于 \(\mathfrak{p}\) 是素理想，存在某个 \(n\) 使得 \(\mathfrak{b}_n \subset \mathfrak{p}\)（第二章，第 1 节，第 1 号，命题 1），从而 \(\mathfrak{b}_n = \mathfrak{p} = \operatorname{Ann}(x_{i(n)})\)，证明完成。

推论。对于每个与分次 A-模 M 相伴的（必然是分次的）素理想 p，存在 \(k \in \Delta\)，使得由分次 A-模 A/p 将次数减去 k 而得到的分次 A-模 \((\mathbf{A}/\mathfrak{p})(k)\)（代数，第二章，第 11 节，第 2 节）同构于 M 的一个分次子模。

  沿用命题 1 证明中的记号。考虑由 A 到 M 的同态 \(a \mapsto ax_{i(n)}\) 取商得到的同态；后者是次数为 \(i(n)\) 的分次同态，因而取商后给出一个次数为 \(i(n)\) 的分次双射同态，把 \(A / p\) 映到 \(M\) 的一个分次子模上。

命题 2. 设 \(\Delta\) 是无挠交换群，A 是 \(\Delta\) 型分次诺特环，M 是 \(\Delta\) 型分次有限生成 A-模。存在一个合成列 \((\mathbf{M}_i)_{0 \leqslant i \leqslant n}\)，其各项都是 M 的分次子模，且对于 \(0 \leqslant i \leqslant n - 1\)，分次模 \(\mathbf{M}_i / \mathbf{M}_{i+1}\) 同构于移位分次模 \((\mathbf{A} / \mathfrak{p}_i)(k_i)\)，其中 \(\mathfrak{p}_i\) 是 A 的分次素理想，\(k_i \in \Delta\)。

只需重走 § 1，第 4 节，定理 1 的论证，这次取 \(\mathfrak{E}\) 为 M 的分次子模所成的集合，其中每个子模都有具备陈述性质的合成列；最后使用命题 1 的推论即可。

